\begin{abstract}
Graph Neural Networks (GNNs) have demonstrated strong performance for
cybersecurity anomaly detection tasks — including network intrusion
classification, botnet identification, and IoT malware analysis — yet two
critical limitations remain unresolved simultaneously: the inability to
handle inherently imprecise network telemetry through fuzzy reasoning, and
the lack of distribution-free, finite-sample coverage guarantees that
compliance-sensitive environments demand.
We present \textbf{FC-GNN} (\textbf{F}uzzy-\textbf{C}onformal
\textbf{G}raph \textbf{N}eural \textbf{N}etwork), the first architecture to
jointly address both weaknesses.
FC-GNN replaces the standard crisp aggregation function with a
\emph{fuzzy message-passing layer} that computes Gaussian membership
functions and Product T-norm rule-firing strengths, yielding
uncertainty-aware node embeddings.
A Sugeno-type fuzzy integral then defines a nonconformity score that feeds
into \emph{Fuzzy Mondrian Conformal Prediction} (FMCP), which partitions
calibration nodes into graph communities via Louvain clustering and
computes per-community quantiles — providing conditional, finite-sample
coverage guarantees of the form
$\Prob(y_v \in \Calpha(v) \mid v \in Q_m) \ge 1-\alpha$.
We prove a coverage theorem for FC-GNN under a fuzzy-extended
permutation-invariance condition, and show that the conditional slack is
$O(|Q_m|^{-1/2})$.
Experiments on seven publicly available cybersecurity benchmarks
— CIC-IDS-2017, UNSW-NB15, NF-BoT-IoT, ISCX-Botnet 2014, CTU-13,
N-BaIoT, and ToN-IoT — confirm that FC-GNN is the only method to
simultaneously achieve valid conditional coverage ($\ge\!90\%$ on $5/7$
datasets at $\alpha\!=\!0.10$) while outperforming the leading Mondrian
CP baseline (RR-GNN) in prediction-set efficiency on three datasets, and
providing human-interpretable IF-THEN rule explanations aligned with
MITRE ATT\&CK categories (mean explanation stability $0.61$).
\keywords{Graph neural networks \and Conformal prediction \and
Fuzzy logic \and Cybersecurity \and Anomaly detection \and
Uncertainty quantification \and Interpretability}
\end{abstract}
